\chapter{Information-Ceiling and Falsifiability Diagnostics}
\label{app:info_ceiling_diagnostics}

\section{Information-Theoretic Ceiling}
\label{sec:info_limit}

Three architecture-independent diagnostics quantify the $K=100$ sensor information limit. The wavelet sparsity diagnostic (Table~\ref{tab:wavelet}) yields Gini $\approx 0.985$ for $u, v$; $99\%$ of energy concentrates in $\sim 328$ of $65\,536$ coefficients ($s/N \approx 0.5\%$). The CS bound \eqref{eq:cs_bound} then gives
\begin{equation}
    M \;\gtrsim\; C_0 \cdot 328 \cdot \log_2(65\,536/328) \;\in\; [2\,500, 5\,000] \quad (C_0\in[1,2]),
\end{equation}
so $K=100$ covers the approximation band ($\sim 196$ coefficients, $94.4\%$ energy) but lies below the CS threshold for mid/high-band detail coefficients. A 2D Fourier basis yields comparable Gini values and the same $\mathcal{O}(s\log_2 N)$ scaling ($M^*=2\,506$ for $C_0=1$); the wavelet-sparsity algorithm is given below.

\begin{table}[!htbp]
    \centering
    \caption{Wavelet sparsity diagnostic on $u, v$ (m/s) and vorticity $\omega$ (1/s) at $t=5$ (db4 2-D wavelet transform). Sparsity threshold $s$ is the number of coefficients carrying $99\%$ of the energy.}
    \label{tab:wavelet}
    \small
    \begin{tabular}{lccc}
    \toprule
    \textbf{Field} & \textbf{Gini coefficient} & \textbf{99\%-energy coefficients} & \textbf{Fraction of total} \\
    \midrule
    $u$ & 0.983 & 326 & 0.50\% \\
    $v$ & 0.985 & 330 & 0.50\% \\
    $\omega$ & 0.942 & 1\,917 & 2.92\% \\
    \bottomrule
    \end{tabular}
\end{table}

\paragraph{Wavelet sparsity algorithm.}
The diagnostic applies, for each $u(t,\cdot)$ and $v(t,\cdot)$: (i) a 5-level db4 2D wavelet decomposition; (ii) coefficients flattened and sorted by absolute value; (iii) a cumulative-energy curve $E_p = \sum_{i\le p} c_i^2 / \sum_i c_i^2$, taking $p^*$ with $E_{p^*}\ge 0.99$; and (iv) the Gini coefficient $G = \frac{2\sum_i i\,c_i}{N\sum_i c_i} - \frac{N+1}{N}$ (with $\{c_i\}$ in ascending $|c|$). Values appear in Table~\ref{tab:wavelet}.

\paragraph{Null-space dimension via sensing SVD.}
A complementary measure counts how many divergence-free Fourier modes the $K=100$ sensors directly constrain. Parameterising the field by its stream function and truncating at $|\mathbf{k}|\le 16$ gives $M = 796$ real degrees of freedom; the sensing matrix $A\in\mathbb{R}^{200\times 796}$ maps these to the $2K=200$ scalar $(u,v)$ measurements. SVD gives $\mathrm{rank}(A) = 200$ (smallest singular value $\approx 497$, well above noise), so the null-space fraction within the $k\le 16$ truncation is $(796-200)/796 = 74.9\%$. In the approximation band ($M=196$, $k\lesssim 8$, the largest $M<2K$ on a 2D radial grid), every divergence-free mode is linearly observable ($\mathrm{rank}(A)=196$). Conditioning still degrades sharply across this band: $\kappa\approx 7$ for $k\le 5$ versus $\kappa\approx 7\times 10^{2}$ for $k\le 8$, so modes near $k\lesssim 8$ are observable but strongly noise-amplifying. The effective reconstructable bandwidth is therefore set by conditioning and noise rather than a hard observability wall, consistent with the achieved effective cutoff $k_{\rm cut}\approx 4.7$ lying inside the $k\lesssim 8$ ceiling.

A complementary spectral-truncation bound assumes perfect reconstruction up to cutoff $k_{\rm cut}$; the residual KE and $\omega$ errors are then architecture-independent (Table~\ref{tab:spectral_truncation}). The $K=100$ sensor Nyquist gives $k_{\max}^{\rm sensor} = \sqrt{K/\pi} \approx 5.64$.

\begin{table}[!htbp]
    \centering
    \caption{Spectral-truncation lower bound: residual KE and vorticity relative error after perfect reconstruction up to cutoff $k_{\rm cut}$.}
    \label{tab:spectral_truncation}
    \small
    \setlength{\tabcolsep}{8pt}
    \begin{tabular}{ccc}
    \toprule
    \textbf{$k_{\rm cut}$} & \textbf{KE loss \%} & \textbf{$\omega$ rel-$L_2$ \%} \\
    \midrule
    4  & $7.77$ & $63.91$ \\
    5  & $4.85$ & $57.18$ \\
    6  & $2.62$ & $51.07$ \\
    8  & $1.05$ & $40.98$ \\
    16 & $0.09$ & $20.96$ \\
    \bottomrule
    \end{tabular}
\end{table}

\paragraph{Ceiling attribution.}
The B3 multi-seed KE error (Table~\ref{tab:multiseed}, $5.71\%$) sits between the $k_{\rm cut}=4$ ($7.77\%$) and $k_{\rm cut}=5$ ($4.85\%$) bounds, mapping to effective cutoff $k_{\rm cut}\approx 4.7$; the vorticity error ($41.79\%$) sits between the $k_{\rm cut}=6$ ($51.07\%$) and $k_{\rm cut}=8$ ($40.98\%$) bounds, mapping to $k_{\rm cut}\approx 7.8$, at the upper edge of the $k\lesssim 8$ linear-observability band identified above. The two metrics bracket the effective bandwidth rather than agreeing on one value, as follows from fitting a single sharp-cutoff model to an energy-weighted and a $k^2$-weighted error in turn. The residual mid/high-band error is consistent with the sensor-information limit at this budget: reducing it requires more sensors rather than changes to architecture, optimizer, or loss (\S\ref{sec:engineering_app}).
